% !TeX program = lualatex
% !TeX encoding = utf8
% !TeX spellcheck = uk_UA
% !BIB program = biber

\documentclass{PySCFBook}


\begin{document}

%%% --------------------------------------------------------
\subsection{Детальная математика проекции и уравнений для амплитуд}
%% --------------------------------------------------------

Начнём с анзаца CC и уравнения Шрёдингера:
\[
|\Psi_{\text{CC}}\rangle = e^{\hat{T}}|\Phi_0\rangle,\qquad
\hat{H}|\Psi_{\text{CC}}\rangle = E_{\text{CC}}|\Psi_{\text{CC}}\rangle.
\]
Домножим слева на \(e^{-\hat{T}}\) и введём эффективный гамильтониан
\[
\bar{H} \equiv e^{-\hat{T}}\hat{H}e^{\hat{T}}.
\]
Тогда имеем операторное равенство
\[
\bar{H}|\Phi_0\rangle = E_{\text{CC}}|\Phi_0\rangle.
\]

Разложение \(\bar{H}\) получается через коммутатный (Бейкер–Кемпбелл–Хаусдорф) ряд:
\[
\bar{H} = \hat{H} + [\hat{H},\hat{T}] + \tfrac{1}{2!}[[\hat{H},\hat{T}],\hat{T}]
+ \tfrac{1}{3!}[[[\hat{H},\hat{T}],\hat{T}],\hat{T}] + \cdots .
\]
Если \(\hat{H}\) содержит только одно- и двухчастичные члены (как обычно),
то этот ряд фактически обрывается на скобках конечного порядка (необходимо не больше четырёхкратных коммутаторов).
Важно: каждый коммутатор породит комбинации членов \(\hat{T}_n\), поэтому в \(\bar H\)
будут члены, линейные, квадратичные, кубические и т.д. по амплитудам \(t\).

Далее проецируем операторное уравнение на основной детерминант та на все
збуджені детермінанти \(\{|\Phi_\mu\rangle\}\). Это даёт:
\[
E_{\text{CC}} = \langle\Phi_0|\bar H|\Phi_0\rangle,
\qquad
0 = \langle\Phi_\mu|\bar H|\Phi_0\rangle \quad\text{для всех } \mu.
\]
Последние уравнения — это искомая система нелинейных уравнений для амплитуд \(t\).

\bigskip

\paragraph{Разложение по компонентам (указание структуры уравнений).}
Для ясности разобьём гамильтониан на одно- и двухчастичные части:
\[
\hat H = \hat F + \hat V,
\]
где \(\hat F\) — фоковский (одночастичный) оператор, \(\hat V\) — двухчастичный.
Подставляя в ряд коммутаторов, получаем структуру:
\[
\bar H = \hat F + \hat V
+ [\hat F+\hat V,\hat T]
+ \tfrac12 [[\hat F+\hat V,\hat T],\hat T] + \cdots
\]
Проекция на \(\langle\Phi_\mu|\) даёт сумму вкладов разной степени по \(t\):
\[
\underbrace{\langle\Phi_\mu|\hat F+\hat V|\Phi_0\rangle}_{\text{(а)}}
+ \underbrace{\langle\Phi_\mu|[\hat F+\hat V,\hat T]|\Phi_0\rangle}_{\text{(б)}}
+ \underbrace{\tfrac12\langle\Phi_\mu|[[\hat F+\hat V,\hat T],\hat T]|\Phi_0\rangle}_{\text{(в)}}
+ \cdots = 0.
\]
Обратите внимание:
\begin{itemize}
  \item (а) — вклад, независимый от \(t\) (обычно нулевой для \(\mu\neq 0\) при HF).
  \item (б) — члены \emph{линейные} по \(t\) (появляются из одного коммутатора).
  \item (в) — члены \emph{квадратичные} по \(t\) (из двойного коммутатора) и т.д.
\end{itemize}
Именно эти нелинейные (квадратичные, кубические...) по \(t\) члены делают систему нелинейной.

\bigskip

\paragraph{Пример: вид уравнений для CCSD (структурно).}
Ограничимся \(\hat T\approx\hat T_1+\hat T_2\). Тогда для одиночных и двойных проекций получаем структуру уравнений
\[
R_i^a \equiv \langle\Phi_i^a|\bar H|\Phi_0\rangle = 0,\qquad
R_{ij}^{ab} \equiv \langle\Phi_{ij}^{ab}|\bar H|\Phi_0\rangle = 0.
\]
Здесь \(R\) — «остатки» (residuals). Структурно эти выражения содержат:
\[
\begin{aligned}
R_i^a &=\; f_i^a \;+\; \sum_{e} f_a^e\, t_i^e \;-\; \sum_{m} f_m^i\, t_m^a
\;+\; \sum_{m,e} \langle am||ie\rangle\, t_m^e \\
&\quad + \sum_{m,e} \langle am||ef\rangle\, t_{im}^{ef} \;+\; \text{квадратичные и высшие члены }(t_1 t_2,\ t_1^2,\ \ldots) = 0,
\\[4pt]
R_{ij}^{ab} &=\; \langle ij||ab\rangle \;+\; \mathcal{L}_{\text{lin}}(t_1,t_2)
\;+\; \mathcal{Q}(t_1,t_2)\;=\;0,
\end{aligned}
\]
где \(f_p^q\) — элементы фока, \(\langle pq||rs\rangle\) — антисиметризованные двумерные интегралы, а \(\mathcal L,\mathcal Q\)
символически обозначают линейную и квадратичную комбинации амплитуд (точные развертки см. в любой специализированной книге/статье по CCSD).
Ключевой момент: в \(R\) есть как линейные, так и квадратичные по \(t\) члены, поэтому уравнения нелинейны.

\bigskip

\paragraph{Почему итерационный метод?}
Нелинейную систему \(R(t)=0\) решают численно. Стандартная схема:
\begin{enumerate}
  \item задать начальное приближение \(t^{(0)}\) (обычно из MP2: \(t^{(0)}_{ij}^{ab}= \langle ij||ab\rangle / (\varepsilon_i+\varepsilon_j-\varepsilon_a-\varepsilon_b)\), \(t^{(0)}_i^a\approx0\));
  \item на шаге \(k\) вычислить остатки \(R(t^{(k)})\);
  \item получить новое приближение \(t^{(k+1)}\) например из аппроксимации линейной части или по формуле типа
  \[
    t^{(k+1)} = t^{(k)} + \Delta t,\qquad \Delta t \approx -\mathbf{D}^{-1} R(t^{(k)}),
  \]
  где \(\mathbf D\) — диагональная аппроксимация джакобиана (в практике используют упрощённый обновляющий оператор, DIIS-ускорение и пр.);
  \item повторять, пока \(\|R\|\) или изменения \(t\) (или изменение \(E_{\text{CC}}\)) не станут меньше заданного порога.
\end{enumerate}
Таким образом итерация — это способ численного решения нелинейной системы для амплитуд.

\bigskip

\paragraph{Стационарность vs минимум.}
Важно подчеркнуть: решение уравнений \(\langle\Phi_\mu|\bar H|\Phi_0\rangle=0\) даёт стационарное состояние относительно амплитуд (остатки = 0), но это \emph{не} результат явной минимизации энергофункционала (энергия CC не является вариационной функцией относительно \(t\)). Поэтому корректнее говорить о достижении стационарности/сходимости, а не о «минимизации» в вариационном смысле.

\bigskip

\paragraph{Итог (кратко).}
\begin{itemize}
  \item Проекционные уравнения получаются непосредственно из уравнения Шрёдингера после преобразования \(e^{-\hat T}\).
  \item Коммутаторный разложение порождает члены разного порядка по \(t\) — это источник нелинейности.
  \item Нелинейная система решается итерационно: начинают с MP2-приближения и повторяют обновления \(t\) до сходимости.
\end{itemize}


\end{document}
